\section{Questão 1: pesos iniciais e funções de ativação}

\subsection*{a. Descrição da configuração experimental}

A rede tem 16 entradas, duas camadas ocultas com 50 neurônios cada e 7 saídas. Esse tamanho está dentro da faixa
explorada na Questão 3 e serve de ponto de partida. As camadas ocultas usam ReLU, que não satura para ativações positivas
e costuma convergir bem com a inicialização usada pela classe. A saída usa softmax, que produz uma distribuição de
probabilidade sobre as classes. A perda é a entropia cruzada categórica, a escolha padrão para esse tipo de saída.

O otimizador é SGD com momento = 0,9 e LR = 0,01, por até 300 épocas. Usar SGD desde o início mantém esta etapa comparável
com a Questão 2. Foram feitos cinco treinamentos com seeds diferentes e, para justificar a ReLU, os mesmos cinco com tanh e
logistic. Só o conjunto de treino é usado.

\subsection*{b. Resultados}

\tabela{q1_seeds}{Cinco treinamentos com ReLU, no conjunto de treino.}
\tabela{q1_ativacoes}{Comparação das funções de ativação (média de cinco seeds).}
\figura{q1_curvas_seeds}{Perda e acurácia de treino com ReLU, por seed, a partir da quinta época.}{0.8}
\figura{q1_curvas_ativacoes}{Perda (em escala logarítmica) e acurácia de treino por função de ativação.}{0.8}

\subsection*{c. Análise e discussão}

Os cinco treinamentos com ReLU são quase coincidentes: a perda final varia de 0,1460 a 0,1502, e a acurácia de 0,9441 a
0,9463. A inicialização afeta só as primeiras épocas, e o desvio final é pequeno (0,0015 na perda, 0,0009 na acurácia).
As curvas seguem caindo devagar na época 300, com oscilações pequenas, sem sinal de instabilidade.

Quanto à ativação, ReLU e tanh chegam a resultados próximos (perda 0,1480 e 0,1520). A logistic demora muito mais: são
239 épocas até a perda de 0,20, contra 12 da ReLU, e fica com a pior perda final (0,1947). Nas primeiras épocas, sua
acurácia fica parada em cerca de 0,26, a proporção da classe majoritária.

Sobre over e underfitting, os dados de treino sozinhos só permitem uma conclusão parcial. Não há underfitting evidente,
pois a acurácia passa de 0,94 em todas as seeds. O overfitting não pode ser avaliado sem um conjunto separado, o que é
feito a partir da próxima questão.

\subsection*{d. Conclusão da etapa}

O modelo aprende os padrões básicos, com acurácia de treino de 0,945. O treinamento é estável em relação aos pesos iniciais,
e a ReLU é adequada: converge tão bem quanto a tanh, é muito mais rápida que a logistic e não apresentou nenhuma
instabilidade nessa configuração. A topologia 50$\times$50 com ReLU segue como base da Questão 2.

\FloatBarrier
